\section{KZG commitments}\label{kzg}

The goal of a \emph{polynomial commitment scheme} is to have the
following API:

\begin{itemize}
\item
  Peggy has a secret polynomial
  \(P(X) \in {\mathbb{F}}_{q}\lbrack X\rbrack\).
\item
  Peggy sends a short ``commitment'' to the polynomial (like a hash).
\item
  This commitment should have the additional property that Peggy should
  be able to ``open'' the commitment at any \(z \in {\mathbb{F}}_{q}\).
  Specifically:

  \begin{itemize}
  \item
    Victor has an input \(z \in {\mathbb{F}}_{q}\) and wants to know
    \(P(z)\).
  \item
    Peggy knows \(P\) so she can compute \(P(z)\); she sends the
    resulting number \(y = P(z)\) to Victor.
  \item
    Peggy can then send a short ``proof'' convincing Victor that \(y\)
    is the correct value, without having to reveal \(P\).
  \end{itemize}
\end{itemize}

The \emph{Kate-Zaverucha-Goldberg (KZG)} commitment scheme is amazingly
efficient because both the commitment and proof lengths are a single
point on \(E\), encodable in 256 bits, no matter how many coefficients
the polynomial has.

\subsection{The setup}

Remember the notation \(\lbrack n\rbrack ≔ n \cdot g \in E\) defined in
{[}armor{]}. To set up the KZG commitment scheme, a trusted party needs
to pick a secret scalar \(s \in {\mathbb{F}}_{q}\) and publish
\[\left\lbrack s^{0} \right\rbrack,\left\lbrack s^{1} \right\rbrack,\ldots,\left\lbrack s^{M} \right\rbrack\]
for some large \(M\), the maximum degree of a polynomial the scheme
needs to support. This means anyone can evaluate
\(\left\lbrack P(s) \right\rbrack\) for any given polynomial \(P\) of
degree up to \(M\). For example,
\[\left\lbrack s^{2} + 8s + 6 \right\rbrack = \left\lbrack s^{2} \right\rbrack + 8\lbrack s\rbrack + 6\lbrack 1\rbrack.\]
Meanwhile, the secret scalar \(s\) is never revealed to anyone.

The setup only needs to be done by a trusted party once for the curve
\(E\). Then anyone in the world can use the resulting sequence for KZG
commitments.

strong(Remark):\\

\subsection{The KZG commitment scheme}

Peggy has a polynomial \(P(X) \in {\mathbb{F}}_{p}\lbrack X\rbrack\). To
commit to it:

strong(Algorithm):\\
``Creating a KZG commitment'':

\begin{enumerate}
\item
  Peggy computes and publishes \(\left\lbrack P(s) \right\rbrack\).
\end{enumerate}

This computation is possible as \(\left\lbrack s^{i} \right\rbrack\) are
globally known.

Now consider an input \(z \in {\mathbb{F}}_{p}\); Victor wants to know
the value of \(P(z)\). If Peggy wishes to convince Victor that
\(P(z) = y\), then:

strong(Algorithm):\\
``Opening a KZG commitment'':

\begin{enumerate}
\item
  Peggy does polynomial division to compute
  \(Q(X) \in {\mathbb{F}}_{q}\lbrack X\rbrack\) such that
  \[P(X) - y = (X - z)Q(X).\]
\item
  Peggy computes and sends Victor \(\left\lbrack Q(s) \right\rbrack\),
  which again she can compute from the globally known
  \(\left\lbrack s^{i} \right\rbrack\).
\item
  Victor verifies by checking
  \[\operatorname{pair}(\left\lbrack Q(s) \right\rbrack,\lbrack s\rbrack - \lbrack z\rbrack) = \operatorname{pair}(\left\lbrack P(s) \right\rbrack - \lbrack y\rbrack,\lbrack 1\rbrack)\]
  \phantomsection\label{kzg-verify}{}

  and accepts if and only if {[}kzg-verify{]} is true.
\end{enumerate}

If Peggy is truthful, then {[}kzg-verify{]} will certainly check out.

If \(y \neq P(z)\), then Peggy cannot do the polynomial long division
described above. So to cheat Victor, she needs to otherwise find an
element
\[\frac{1}{s - x}\left( \left\lbrack P(s) \right\rbrack - \lbrack y\rbrack \right) \in E.\]
Since \(s\) is a secret nobody knows, there is not any known way to do
this.

\subsection{Multi-openings}\label{multi-openings}

To reveal \(P\) at a single value \(z\), we did polynomial division to
divide \(P(X)\) by \(X - z\). But there is no reason we have to restrict
ourselves to linear polynomials; this would work equally well with
higher-degree polynomials, while still using only a single 256-bit curve
point for the proof.

For example, suppose Peggy wanted to prove that \(P(1) = 100\),
\(P(2) = 400\), \ldots, \(P(9) = 8100\). (We chose these numbers so that
\(P(X) = 100X^{2}\) for \(X = 1,\ldots,9\).)

Evaluating a polynomial at \(1,2,\ldots,9\) is essentially the same as
dividing by \((X - 1)(X - 2)\ldots(X - 9)\) and taking the remainder. In
other words, if Peggy does a polynomial long division, she will find
that
\[P(X) = Q(X)\left( (X - 1)(X - 2)\ldots(X - 9) \right) + 100X^{2.}\]
Then Peggy sends \(\left\lbrack Q(s) \right\rbrack\) as her proof, and
the verification equation is that \[\begin{aligned}
 & \operatorname{pair}(\left\lbrack Q(s) \right\rbrack,\left\lbrack (s - 1)(s - 2)\ldots(s - 9) \right\rbrack) \\
 & = \operatorname{pair}(\left\lbrack P(s) \right\rbrack - 100\left\lbrack s^{2} \right\rbrack,\lbrack 1\rbrack).
\end{aligned}\]

The full generality just replaces the \(100X^{2}\) with the polynomial
obtained from Lagrange interpolation\footnote{https://en.wikipedia.org/wiki/Lagrange\_polynomial}
(there is a unique such polynomial \(f\) of degree \(n - 1\)). To spell
this out, suppose Peggy wishes to prove to Victor that
\(P\left( z_{i} \right) = y_{i}\) for \(1 \leq i \leq n\).

strong(Algorithm):\\
Opening a KZG commitment at \(n\) values:

\begin{enumerate}
\item
  By Lagrange interpolation, both parties agree on a polynomial \(f(X)\)
  such that \(f\left( z_{i} \right) = y_{i}\).
\item
  Peggy does polynomial long division to get \(Q(X)\) such that
  \[P(X) - f(X) = \left( X - z_{1} \right)\left( X - z_{2} \right)\ldots\left( X - z_{n} \right) \cdot Q(X).\]
\item
  Peggy sends the single element \(\left\lbrack Q(s) \right\rbrack\) as
  her proof.
\item
  Victor verifies \[\begin{aligned}
   & \operatorname{pair}(\left\lbrack Q(s) \right\rbrack,\left\lbrack \left( s - z_{1} \right)\left( s - z_{2} \right)\ldots\left( s - z_{n} \right) \right\rbrack) \\
   & = \operatorname{pair}(\left\lbrack P(s) \right\rbrack - \left\lbrack f(s) \right\rbrack,\lbrack 1\rbrack).
  \end{aligned}\]
\end{enumerate}

So one can even open the polynomial \(P\) at \(1000\) points with a
single 256-bit proof. The verification runtime is a single pairing plus
however long it takes to compute the Lagrange interpolation \(f\).

\subsection{Root check}

To make PLONK work, we are going to need a small variant of the
multi-opening protocol for KZG commitments
(\hyperref[multi-openings]{{[}multi-openings{]}}), which we call
\emph{root check} (not a standard name). Here is the problem statement:

strong(Problem):\\

Peggy just needs to show that \(P_{1} - P_{2}\) is divisible by
\(Z(X) ≔ \prod_{z \in S}(X - z)\). This can be done by committing the
quotient \[H(X) ≔ \frac{P_{1}(X) - P_{2}(X)}{Z(X)}.\] Victor then gives
a random challenge \(\lambda \in {\mathbb{F}}_{q}\), and then Peggy
opens \(\operatorname{Com}(P_{1})\), \(\operatorname{Com}(P_{2})\), and
\(\operatorname{Com}(H)\) at \(\lambda\).

But we can actually do this more generally with \emph{any} polynomial
expression \(F\) in place of \(P_{1} - P_{2}\), as long as Peggy has a
way to prove the values of \(F\) are correct. As an artificial example,
if Peggy has sent Victor \(\operatorname{Com}(P_{1})\) through
\(\operatorname{Com}(P_{6})\), and wants to show that
\[P_{1}(42) + P_{2}(42)P_{3}(42)^{4} + P_{4}(42)P_{5}(42)P_{6}(42) = 1337,\]
she could define
\[F(X) ≔ P_{1}(X) + P_{2}(X)P_{3}(X)^{4} + P_{4}(X)P_{5}(X)P_{6}(X) - 1337\]
and run the same protocol with this \(F\). This means she does not have
to reveal any \(P_{i}(42)\), which is great!

To be fully explicit, here is the algorithm:

strong(Algorithm):\\
Root check: Assume that \(F\) is a polynomial for which Peggy can
establish the value of \(F\) at any point in \({\mathbb{F}}_{q}\). Peggy
wants to convince Victor that \(F\) vanishes on a given finite set
\(S \subseteq {\mathbb{F}}_{q}\).

\begin{enumerate}
\item
  If she has not already done so, Peggy sends to Victor a commitment
  \(\operatorname{Com}(F)\) to \(F\).\footnote{In fact, it is enough for
    Peggy to have some way to prove to Victor the values of \(F\).

    So for example, if \(F\) is a product of two polynomials
    \(F = F_{1}F_{2}\), and Peggy has already sent commitments to
    \(F_{1}\) and \(F_{2}\), then there is no need for Peggy to commit
    to \(F\).

    Instead, in Step 5 below, Peggy opens \(\operatorname{Com}(F_{1})\)
    and \(\operatorname{Com}(F_{2})\) at \(\lambda\), and that proves to
    Victor the value of \(F(\lambda) = F_{1}(\lambda)F_{2}(\lambda)\).}
\item
  Both parties compute the polynomial
  \[Z(X) ≔ \prod_{z \in S}(X - z) \in {\mathbb{F}}_{q}\lbrack X\rbrack.\]
\item
  Peggy does polynomial long division to compute
  \(H(X) = \frac{F(X)}{Z(X)}\).
\item
  Peggy sends \(\operatorname{Com}(H)\).
\item
  Victor picks a random challenge \(\lambda \in {\mathbb{F}}_{q}\) and
  asks Peggy to open \(\operatorname{Com}(H)\) at \(\lambda\), as well
  as the value of \(F\) at \(\lambda\).
\item
  Victor verifies \(F(\lambda) = Z(\lambda)H(\lambda)\).
\end{enumerate}

\phantomsection\label{root-check}{}
